% id: 0428
% book: RPM 중학 3-1 · 04 다항식의 곱셈
% section: 유형 07 곱셈 공식 종합
% figure: crop:fig-0428.png
% answer: ③
% answer_source: 답지
% variant_level: 0 (원본 전사)
% 이 파일은 build.mjs 가 items.json 에서 만든다. 고칠 때는 items.json 을 고친다.
\smash{\raisebox{1.5mm}{\dmkichul{RPM 중학 3-1 0428번 · 61쪽}}}
\noindent\begin{minipage}[t]{0.58\linewidth}\vspace{0pt}\raggedright 오른쪽 그림의 세 직사각형의 넓이 $P$, $Q$, $R$에 대하여 $P+Q=P+R$일 때, 다음 중 설명할 수 있는 곱셈 공식으로 가장 적당한 것은?\end{minipage}\hfill\begin{minipage}[t]{0.4\linewidth}\vspace{0pt}\centering \includegraphics[width=81.7pt]{figures/fig-0428.png}\end{minipage}\par
\begin{choicesv}{$(a+b)^2=a^2+2ab+b^2$}{$(a-b)^2=a^2-2ab+b^2$}{$(a+b)(a-b)=a^2-b^2$}{$(x+a)(x+b)=x^2+(a+b)x+ab$}{$(ax+b)(cx+d)=acx^2+(ad+bc)x+bd$}\end{choicesv}
